% s23 Open problems
\section{Open problems}

Throughout, $K$ denotes King's printed state of \cref{thm:s13:main}: the Neumann--Dirichlet cube limit on King's RG cube at block factor $5$, fixed photon mass $\lmu>0$, with an admissible scheme $a$ in the Wick-square sector; $S$ denotes a torus limit state of the authors' earlier paper \cite{SymPhase} (block factor $2$, torus counterterm scheme with reference photon mass $m$, $0<\lambda,e^2<0.1m$, $\lmu\in(0,m]$). No identification between the two is proved. ``Symmetric phase'' is, for both, the expectation that the states lie above the surface predicted by one-loop perturbation theory; it is neither a hypothesis nor a conclusion of any theorem here, and \cref{thm:s13:main} holds on all of King's coupling range (\cref{lem:s14:CW}), the surface included. The problems below are extensions of the main theorem, not conditions on it.

\subsection{Relation to the torus states}

Three hypotheses bear on the relation of $K$ to $S$: $(\mathrm{FAM}_L)$, King's cube construction exists on the block-factor-$2$ family in the torus counterterm scheme; $(\mathrm{KI}_L)$, the fixed-torus continuum limit exists in that family; $(\mathrm{CE})$, a convergent unit-scale cluster expansion with sources on the cubic tori, their tilted copies and their cuts along King's cubes, uniform in the lattice spacing. Under $(\mathrm{FAM}_L)$ and a sign or decay hypothesis on the difference of the two counterterms, every torus state dominates the King state at a matched bare mass on the observables $e^{\mathrm iF(f)+\Wick{k}}$, $k\ge0$, a one-sided inequality carrying no information about rotations. Under $(\mathrm{CE})$ at a point of a massive window ($\lambda m^2\ge K_0e^6$, $K_0\ge1/(18\pi^2)$, small $e^2,\lambda$) every torus limit state equals the cube limit of the same family and scheme in every order of limits; with $(\mathrm{FAM}_L)$ or $(\mathrm{KI}_L)$ it is unique, exponentially clustering and ergodic; with King's cube limits on both families it equals $K$ at an explicitly shifted bare mass. Whether $(\mathrm{CE})$ holds anywhere at $e>0$ is open. Two obstacles are structural: the torus paper computes its mass counterterm with the photon propagator frozen at $m$ while King's free-boundary counterterm uses the model's own $\lmu$, a difference of order $e^2|m-\lmu|$, not small (\cref{rem:s14:torus-window}); and King's construction needs block factor $5$ for the rotation by $\arctan(3/4)$, while no lattice family admits both a dyadic structure and a rotation of infinite order. On the one-loop surface $\lambda m^2=e^6/(72\pi^2)(1+O(e^2/m))$ the two constructions are expected to select different phases (periodic tori the broken one, Dirichlet faces the symmetric one); this is a heuristic, not a theorem, and no statement of this paper uses the phase.

\subsection{Uniqueness and ergodicity of \texorpdfstring{$K$}{K}}

Uniqueness of King's cube limit is part of the construction (monotonicity in the cube), but uniqueness of the infinite-volume state, in the sense of King's expected identity between cube limits and periodic-box limits, is not proved, and neither is ergodicity of $K$ under translations: \cref{thm:s13:main} gives clustering only on the ergodic components. The usual route, a cluster expansion with a mass gap, is $(\mathrm{CE})$. This is distinct from \UNIQ\ (\cref{hyp:s12:UNIQ}), the independence of a LIM state from its continuum subsequence, which follows from \cref{hyp:s03:EC} and from \cref{hyp:s03:RC} (\cref{prop:s12:ID}, \cref{cor:s04:G2}). Cubic symmetry of $K$ is a theorem (\cref{thm:s06:S1K}): it needs only the uniqueness of the monotone cube limit and the reflection sandwich (\cref{rem:s06:K4}).

\subsection{Non-triviality}

For the torus states the truncated correlation of $F(g)^2$ with the Wick square is non-positive, and strictly negative at small couplings under the two hypotheses of \cite[thm:NT-cond]{SymPhase}: the one-loop statement (C) on the continuum coefficient and a remainder bound on it, which is unproved and, as printed, is expected to fail below the first-order surface. The analogous statements for King's cube, and a non-vanishing two-point function, are left to later work. For $K$ nothing is proved here beyond the sign inequalities that hold on every lattice. Whether $K$ is non-Gaussian, or even not decoupled, is open.

\subsection{The limit of vanishing photon mass}

Every constant of the Wick leg of \cref{thm:s13:main} depends on $\lmu$ and is not uniform in it (\cref{rem:s13:mu0}); the absolute bounds of the curvature leg are. The limit $\lmu\downarrow0$, taken along a frozen family (counterterm and scheme frozen at a reference mass, so that the members are monotone in $\mu$ on the cone by the correlation inequalities), is left to later work: in the curvature sector it needs the existence theorem for every member and nothing else; in the Wick sector it needs a cell bound uniform in the photon mass, together with the torus inputs at the reference mass and the transfer of \cref{thm:s13:TWK} for the members on King's convention; the identification of the members with King's models at shifted scalar masses, masslessness proper (the vanishing of the photon mass in the limit theory) and the two-point function are further questions there. None of these enters the present paper.

\subsection{The optional extension: \texorpdfstring{\Wtwo}{(W2)} at arbitrary amplitudes, Form (A), and the Euclidean clause of the Wick sector}

Form~(B) of \cref{thm:s13:main} lives on the tube $\Dtube$: the only positive-source bound it uses is the one proved on the cone $\sup q\le s_1<\sigma_1(M_1)$ (\cref{cor:s11:W2prime}), and the only existence statement is \cref{hyp:s03:W1}, whose positive scalar part is below the same threshold. Above the threshold, the existence of the fixed-cube limits and the uniform positive-source bound are the hypothesis \Wtwo\ (\cref{hyp:s03:W2}), which is open.

\begin{theorem}[Form (A): the Wick leg at arbitrary amplitudes]\label{thm:s23:FormA}
Let the couplings satisfy \Tthree, and let $a$ be an admissible scheme satisfying \Wtwo\ for the comparison convention $\NDa$. Then for every finite family $(f_1,\dots,f_n;h_1,\dots,h_p)$ of real $2$-forms and real scalars the full object $Z^a_\infty(\zeta)=\lim_J\lim_kZ^a_{k,J}(\sum_iz_if_i,\sum_jw_jh_j)$ exists for every $\zeta=(z,w)\in\C^{n+p}$, locally uniformly, $Z^a_\infty$ is entire, and $Z^a_\infty=\Ginf$ on $\Dtube$. Its Schwinger functions are those of Form~(B), and:
\begin{enumerate}[label=\textup{(\roman*)},leftmargin=*,nosep]
\item \OSz\ as in Form~(B)(i): the tube bound on $\Dtube$ and the bounds on the $S_{n,p}$; no quantitative bound on $Z^a_\infty$ off the tube is asserted.
\item \OSo: $Z^a_\infty(\tau_bR\cdot\zeta)=Z^a_\infty(\zeta)$ for every $\zeta\in\C^{n+p}$, $b\in\R^3$, $R\in\Oh$ ($F$ a $2$-form, $W^a$ a scalar).
\item \OSt\ in matrix form about every coordinate plane for exponentials with arbitrary complex coefficients (\cref{thm:s23:OS2W}), and in polynomial form (\cref{prop:s08:poly}).
\item \OSth\ as in Form~(B)(iv).
\item[\textup{(v),(vi)}] As in Form~(B), the characteristic functional $(u,v)\mapsto Z^a_\infty(\mathrm iu,\mathrm iv)$ being the restriction of the entire $Z^a_\infty$ to the imaginary axes.
\end{enumerate}
\emph{Euclidean clause.} If in addition \RW\ holds for the comparison family, $Z^a_\infty$ is invariant under every Euclidean motion, its mixed Schwinger functions are $\Euc$-covariant, and $\nu$-almost every ergodic component of the measure of (v) is $\Euc$-invariant and reflection positive about every plane (no vanishing of odd moments is asserted for a component). This is the only Euclidean clause of the Wick sector.

\emph{On King's convention.} The same statements for the K-family object $Z^{a,K}_\infty$ under \Wtwo\ read for $\NDK$ are the content of \FAK\ (\cref{hyp:s03:FAK}); nothing is derived for it from the comparison-family statements, since \Wtwo\ on $\NDK$ does not give \Wtwo\ for the comparison family (\cref{cor:s04:G2}(iv)).
\end{theorem}

\begin{proof}[Indications]
Existence at all complex coefficients, locally uniform convergence and entireness are \cref{lem:s08:W1} under \Wtwo\ (version (A$_\infty$) of \cref{lem:s08:claim4}); $Z^a_\infty=\Ginf$ on $\Dtube$ as the same $J$-limit of the same $Z^a_J$ (\cref{thm:s11:W}). (i) is Form~(B)(i). (ii): Form~(B)(ii) is the identity on $\Dtube$, an open neighbourhood of the origin of $\C^{n+p}$, and both sides are entire. (iii) is \cref{thm:s23:OS2W} with \cref{prop:s08:poly}; (iv) is \cref{prop:s08:OS3} for $Z^a_\infty$; (v),(vi): the measure of Form~(B) is built from $\Ginf$ on the imaginary axes, which lie in $\Dtube$. The Euclidean clause is \RW\ itself; the component clause follows through the ergodic decomposition of Form~(B)(vi) as in the curvature sector (\cref{thm:s09:S2K}(c)).
\end{proof}

For a LIM state of a continuum subsequence (\cref{thm:s12:LIM}) the corresponding Form~(A) requires both clauses of \Wtwo: the existence of the limits at every real coefficient vector and the uniform bound at every amplitude.

\begin{theorem}[reflection positivity at arbitrary complex coefficients]\label{thm:s23:OS2W}
Let $a$ be an admissible scheme satisfying \Wtwo, let $\nu\in\{1,2,3\}$, $c\in\R$, and let $\zeta_1,\dots,\zeta_N$ be finite families of real sources $f_{il}\in C^\infty_c(\R^3;\Lambda^2)$, $h_{il}\in C^\infty_c(\R^3)$ supported in $\{x_\nu>c\}$, with arbitrary complex coefficients $z_{il},w_{il}$. Then
\[
\sum_{i,j=1}^N\bar\gamma_i\gamma_j\,Z^a_\infty\big(\theta\zeta_i^*\sqcup\zeta_j\big)\ \ge\ 0\qquad\text{for all }\gamma\in\C^N .
\]
\emph{Refinement.} Let $B_1$, $B''$, the vertex sources $q_v$ ($v\in\vertex(B'')$) and $q_\zeta$ be those of \cref{def:s08:enlarged} for the plane $\{x_\nu=c\}$. (a) If $c\in D$, the proof uses \Wtwo(b) only at the finitely many sources $q_v$, and it suffices to assume it at the single source $(1+\eta)q_\zeta$ for any one $\eta>0$ (\cref{rem:s08:oneinstance}, one direction only); all other steps use \Wtwo(a) and the inequalities of \cref{lem:s02:BFS}, and the tube machinery is not needed. (b) If $c\notin D$, with $c'\in D$ and $b\defeq(c'-c)e_\nu$: the finitely many instances $q_v$ and $\tau_bq_v$, $v\in\vertex(B'')$, together with the tube inputs \Itwo, \Ithree\ of \cref{thm:s08:OS2WT}. Such a plane needs a non-lattice translation, which finitely many instances of \Wtwo(b) do not supply; it is carried by the translation invariance of the tube object (\cref{prop:s08:T}(b)).
\end{theorem}

\begin{proof}
Let $c\in D$ and let $\cF$, $\Phi_k$, $\zeta^*_{ij}$, $\cT_1$, $T_\cF$, $U$ be as in \cref{def:s08:enlarged}. By \cref{lem:s08:W1}, under \Wtwo(a) and \Wtwo(b) at the $q_v$, $\{Z^a_J\}_{J\ge J_0}$ is uniformly bounded on $\cT_1$, $Z^a_J\to Z^a_\infty$ locally uniformly on $\cT_1$, $Z^a_\infty$ is holomorphic there and $Z^a_J\le Z^a_\infty$ on $U\cap\cT_1$. By \cref{lem:s08:claim4}(A), $\limsup_k\sup_{\cT_1}|\Phi_k|\le C$ with $C$ independent of $J\ge J_0$, and at fixed $J$ the $\Phi_k$ are locally bounded on $\C^m$ (\Wtwo(a) only). Montel's theorem gives a subsequence along which $\Phi_k\to\Phi^J$ locally uniformly, $\Phi^J$ entire; the matrices $(\Phi_k(\zeta^*_{ij}))_{ij}$ are positive semidefinite for large $k$ (\cref{lem:s07:RPbox}(b)), hence so is their limit, and $\sup_{\cT_1}|\Phi^J|\le C$ uniformly in $J$. On $U$ every lattice value is real and the scheme multipliers of the two states agree, so the $\cI$ halves of \cref{lem:s07:L,lem:s07:U} give $Z^a_{k,J}\le\Phi_k\le Z^a_{k,J+2n}$, hence $Z^a_J\le\Phi^J\le Z^a_{J+2n}$ on $U$. On $U\cap\cT_1$, $Z^a_J\uparrow Z^a_\infty$ (\cref{lem:s08:W1}), so $\Phi^J\to Z^a_\infty$ there; by the uniform bound on $\cT_1$, Montel's theorem and the uniqueness set $U\cap\cT_1$ (\cref{def:s08:enlarged}), $\Phi^J\to Z^a_\infty$ locally uniformly on the connected $\cT_1$, in particular at the $\zeta^*_{ij}$, and positivity passes to the limit. The only instances of \Wtwo(b) used are the $q_v$ of \cref{lem:s08:W1,lem:s08:claim4}; \cref{rem:s08:oneinstance} replaces them by $(1+\eta)q_\zeta$.

Let $c\notin D$, $c'\in D$, $b\defeq(c'-c)e_\nu$. The $\tau_b\zeta_i$ are supported in $\{x_\nu>c'\}$ and $\theta_{c'}(\tau_b\zeta)^*=\tau_b(\theta_c\zeta^*)$, so the display at $c$ follows from the display at $c'$ for the translated families once $Z^a_\infty$ is translation invariant at the points in question. The enlarged family of the translated problem at $c'$ is $\tau_b\cF$, with the same coefficient space, matrix points, $B_1$, $B''$, and vertex sources $\tau_bq_v$; the argument above for $\tau_b\cF$ gives the positivity of $(Z^a_\infty(\tau_b\cdot\zeta^*_{ij}))_{ij}$, and $\zeta\mapsto Z^a_\infty(\tau_b\cdot\zeta)$, $\zeta\mapsto Z^a_\infty(\zeta)$ are holomorphic on $\cT_1$ by \cref{lem:s08:W1}. On the non-empty open set $\cT_1\cap T_\cF$ both equal the tube limits $Z^a_{\infty,T}(\tau_b\cdot\zeta)$, $Z^a_{\infty,T}(\zeta)$, which agree by \cref{prop:s08:T}(b); by the identity theorem they agree on $\cT_1$, in particular at the $\zeta^*_{ij}$.
\end{proof}

\begin{hypothesis}[\Bminus{u}: stability of the shifted torus model]\label{hyp:s23:Bminus}
Let $u>0$. In the free convention and the scheme $a$,
\[
\sup_{k\ge0}\ \limsup_{J\to\infty}\ \alpha^a_{k,T_J}(u)<\infty,\qquad\alpha^a_{k,T}(u)=|T|^{-1}\ln N^a_{k,T}(0,u\mathbf 1),
\]
the free-energy density of the torus functional at the constant scalar source $u$; that is, \Itwo\ of \cref{def:s10:I2} at the amplitude $s_1=u$.
\end{hypothesis}

For $u$ below the threshold $s_*$ of \cref{lem:s14:CW}, \Bminus{u} is \cref{prop:s10:U} (or \cref{prop:s10:Upm}) at an auxiliary mass $M_1$ with $\sigma_1(M_1)>u$, so Form~(B), whose constants are built from \Itwo\ at one such amplitude, already contains every amplitude at which the positive-source bound is known. For $u\ge s_*$ the hypothesis is the exact obstruction. By \cref{lem:s10:M}, exact at fixed $(k,T)$, $\alpha^z_{k,T}(u)$ is, up to explicit $c$-numbers, $|T|^{-1}$ times the difference of the free energies of King's torus model with bare scalar coefficient $\alpha^{\rm tor}_{k,T}(m)-2u$ and with $\alpha^{\rm tor}_{k,T}(m)$, both with the counterterms of the mass $m$: the second is \cref{hyp:s03:T1} at $m$, the first the ultraviolet stability of the shifted model. For $u<s_*$ the shifted coefficient is that of King's model at a positive mass $M_1$; for $u\ge s_*$ it lies, asymptotically in $k$, below $\inf_{M>0}\alpha^\infty_k(M)$, and no positive scalar mass realises it. The printed proofs of \cite{Ba1,Ba2} expand about the Gaussian at the positive scalar mass $m$ and use its mass term in the small-field bounds; no printed theorem covers a bare coefficient below every positive-mass value. Conversely \Bminus{u} is sufficient for \Wtwo(b): \cref{prop:s11:W1prime} at $s_1=u$, whose steps do not depend on the amplitude, gives $\ln N^a_{k,X}(0,q)\le(C_U/u)\norm q_{1,k}$ for $0\le q\le u$ on every Neumann--Dirichlet box, hence \Wtwo(b) in $L^1$ form on $\sup q\le u$. It is not asserted that \Wtwo\ implies \Bminus{u}: the inequalities available bound a constant-source density from below by local sources, never from above, so \Bminus{u} is sufficient and not shown necessary. Also open: an $L^1$ bound for the negative part of the Wick source on cubes, as on the torus paper's tori (the tube bound \eqref{eq:s11:W} carries the weighted norm $\wnorm{k_-}$).

\subsection{Rotation invariance of the Wick sector}

\Cref{thm:s13:main} gives invariance under $\Euc$ for the curvature sector; for the Wick sector it gives invariance under $\Cub$ on the tube and no Euclidean clause. Euclidean invariance of the joint field is the hypothesis \RW\ of the optional extension (\cref{thm:s23:FormA}), and on King's convention \FAK; whether the heavy-shell comparison of the rotation part extends to the Wick square, on the tube or at arbitrary amplitudes, is open. Cubic symmetry is a theorem in every case.

\subsection{The coupling condition}

The sign clause \Tthree\ is proved on the explicit regions of \cref{prop:s14:T3} (at $\lmu=m$ on a region containing $0\le\lambda,e^2\le0.1m$; at every $\lmu\in(0,m]$ on explicit boxes; exactly for $\lambda<\frac\pi4(1+\sqrt5)\,m$ at $e=0$), and the Wick leg is stated there. Outside those regions the exact position of the failure set $E^<_W$ is not determined; the two-sided clause that $m$ is not a global minimiser of $\alpha_\infty$, with the auxiliary mass on either side of $m$ (\cref{prop:s10:Upm}, \cref{lem:s14:CW}(b)), would enlarge the region. The three hypotheses \Bcell, \SCH, \UNIQ\ of the statements on continuum subsequences (\cref{hyp:s12:Bcell,hyp:s12:SCH,hyp:s12:UNIQ}) are proved for the objects of the main theorem (\cref{prop:s12:ID}) and are open only for LIM states of schemes that are not admissible. The mass clause of \cite[Prop.~3.1]{Ba2}, a cap on the dimensionless scalar mass at the stopping scales where that proposition is used, belongs to the statement of the proposition, which is cited under it, and imposes nothing on the bare mass (\cref{rem:s14:BALm}).
